% !TEX program = lualatex
% ---------------------------------------------------------------------------------------------
%  Readback note 11: The volume of the chain of theta lifts as a product of lift volumes
%
%  Build:  python3 tools/rb.py note <slug>
%
%  One unit of a Lean artifact, rewritten as ordinary mathematics. The prose is hand-written. Lean names,
%  line ranges and counts come from data.tex; quoted Lean statements are copied from the artifact by the build.
%
%  Rules of the text: standard terminology; every symbol is explained where it first appears and listed again
%  in the notation at the end; no letter has two meanings; a proof is complete or says what it takes for granted;
%  Lean names appear only in the LEAN lines and in the concordance.
%
%  The signature rule (the build enforces it): a function, relation or predicate that the artifact defines is
%  written only through its notation macro, which the build makes from the spec (`notation`) with one slot for
%  each explicit argument of the Lean definition. No argument is fixed and dropped, also not "for brevity" and
%  not for one section. Every argument and every hypothesis of a stated declaration stands inside its
%  statement, or in a Setting that the statement cites by number; only implicit and instance arguments (the
%  ambient objects) may stand in the conventions.
% ---------------------------------------------------------------------------------------------
\documentclass[11pt,a4paper]{article}
\ifdefined\pdfvariable\pdfvariable suppressoptionalinfo \numexpr 1+2+4+8\relax\fi
\usepackage[a4paper,hmargin=35mm,top=22mm,bottom=27mm,footskip=12mm,nohead]{geometry}

\usepackage{fontspec}
\directlua{
  luaotfload.add_fallback("gfallback", {"STIXTwoMath-Regular.otf:mode=node;", "DejaVuSans.ttf:mode=node;"})
  luaotfload.add_fallback("gmonofallback", {"DejaVuSansMono.ttf:mode=node;", "STIXTwoMath-Regular.otf:mode=node;"})
}
\setmainfont{NewCM10}[Extension=.otf,UprightFont=*-Book,ItalicFont=*-BookItalic,
  BoldFont=*-Bold,BoldItalicFont=*-BoldItalic,SmallCapsFont=*-Book,SmallCapsFeatures={Letters=SmallCaps},
  RawFeature={fallback=gfallback}]
\setsansfont{NewCMSans10}[Extension=.otf,UprightFont=*-Book,ItalicFont=*-BookOblique,
  BoldFont=*-Bold,BoldItalicFont=*-BoldOblique,RawFeature={fallback=gfallback}]
\setmonofont{JetBrainsMono}[Extension=.otf,UprightFont=*-Regular,ItalicFont=*-Italic,
  BoldFont=*-SemiBold,BoldItalicFont=*-SemiBoldItalic,Scale=MatchLowercase,
  RawFeature={-calt;-liga;fallback=gmonofallback}]
\usepackage{amsmath}
\usepackage{amsthm}
\usepackage[math-style=TeX]{unicode-math}
\setmathfont{NewCMMath-Book.otf}
\emergencystretch=2.5em\relax

\usepackage{xcolor}
\usepackage{tikz}
\usetikzlibrary{calc,arrows.meta,positioning,backgrounds}
\usepackage[protrusion=true,expansion=true]{microtype}
\usepackage{enumitem}
\usepackage{needspace}
\usepackage{booktabs}
\usepackage{array}
\usepackage{longtable}
\usepackage{fancyhdr}
\usepackage{lastpage}
\usepackage{multicol}
\usepackage[hidelinks,pdfusetitle]{hyperref}

\definecolor{ink}{HTML}{0F172A}
\definecolor{inksoft}{HTML}{475569}
\definecolor{inkfaint}{HTML}{94A3B8}
\definecolor{rulegrey}{HTML}{CBD5E1}
\definecolor{arith}{HTML}{C2410C}
\definecolor{est}{HTML}{7C3AED}
\definecolor{ana}{HTML}{0891B2}
\definecolor{ok}{HTML}{047857}
\definecolor{leancol}{HTML}{3730A3}

\title{The volume of the chain of theta lifts as a product of lift volumes}
\author{grh-compact}
\date{}

% ------------------------------------------------------------------------------------- macros
\newcommand{\leanfont}{\ttfamily\fontsize{8.6}{10.5}\selectfont\color{leancol}}
\newcommand{\lean}[1]{\csname LN@#1\endcsname}
\newcommand{\fillin}[1]{{\color{arith}\itshape[#1]}}  % a place of the template that is not written yet; the build reports every one that is left
\newcommand{\leanquote}[1]{\csname LQ@#1\endcsname}
% generated by rb.py note; do not edit
\newcommand{\ConcDecls}{8}
\newcommand{\Pin}{72fead10}
\newcommand{\PinText}{that of commit {\leanfont 72fead10}}
\newcommand{\ArtifactName}{grh-compact}
\expandafter\newcommand\csname LN@normDetSq\endcsname{\href{https://github.com/lsernau/grh-lean/blob/72fead102f663d5c7a407964a13dc0cf9d977207/grh-compact/src/New/Rows/RowVolume.lean\#L15-L19}{{\leanfont normDet\_sq\_eq\_pow\_mul\_det}}}
\expandafter\newcommand\csname LN@liftVolume\endcsname{\href{https://github.com/lsernau/grh-lean/blob/72fead102f663d5c7a407964a13dc0cf9d977207/grh-compact/src/New/Rows/RowVolume.lean\#L23-L23}{{\leanfont Data.\allowbreak{}liftVolume}}}
\expandafter\newcommand\csname LQ@liftVolume\endcsname{\leanquoteline{0}{noncomputable def Data.liftVolume (ϑ : D.MFoldTwist) j := ((D.goodLift (ϑ.pow j) j).normDet : ℂ) \textasciicircum{} 2}}
\expandafter\newcommand\csname LN@chainTheta\endcsname{\href{https://github.com/lsernau/grh-lean/blob/72fead102f663d5c7a407964a13dc0cf9d977207/grh-compact/src/New/Rows/RowVolume.lean\#L27-L27}{{\leanfont Data.\allowbreak{}Θ}}}
\expandafter\newcommand\csname LQ@chainTheta\endcsname{\leanquoteline{0}{noncomputable abbrev Data.Θ (ϑ : D.MFoldTwist) := D.iotaAtLevel ϑ ∘ₗ D.rows.subtype}}
\expandafter\newcommand\csname LN@chainVolume\endcsname{\href{https://github.com/lsernau/grh-lean/blob/72fead102f663d5c7a407964a13dc0cf9d977207/grh-compact/src/New/Rows/RowVolume.lean\#L30-L30}{{\leanfont Data.\allowbreak{}𝔷}}}
\expandafter\newcommand\csname LQ@chainVolume\endcsname{\leanquoteline{0}{noncomputable def Data.𝔷 ϑ := ((D.Θ ϑ).normDet : ℂ) \textasciicircum{} 2}}
\expandafter\newcommand\csname LN@volumeProduct\endcsname{\href{https://github.com/lsernau/grh-lean/blob/72fead102f663d5c7a407964a13dc0cf9d977207/grh-compact/src/New/Rows/RowVolume.lean\#L37-L58}{{\leanfont Data.\allowbreak{}volume\_iota\_eq\_prod\_liftVolume}}}
\expandafter\newcommand\csname LQ@volumeProduct\endcsname{\leanquoteline{0}{lemma Data.volume\_iota\_eq\_prod\_liftVolume : ∃ c, D.CostsSquares (D.m * D.m) c ∧ ∀ ϑ, D.𝔷 ϑ ≠ 0 ∧}\leanquoteline{2}{D.𝔷 ϑ = ‖c ϑ‖ \textasciicircum{} D.rowDim * ∏ j ∈ Icc 1 D.m, D.liftVolume ϑ j}}
\expandafter\newcommand\csname LN@rangeTupleIdem\endcsname{\href{https://github.com/lsernau/grh-lean/blob/72fead102f663d5c7a407964a13dc0cf9d977207/grh-compact/src/New/Rows/RowVolume.lean\#L62-L65}{{\leanfont Data.\allowbreak{}range\_tupleIdem\_le}}}
\expandafter\newcommand\csname LN@groupAlgActionAuxLift\endcsname{\href{https://github.com/lsernau/grh-lean/blob/72fead102f663d5c7a407964a13dc0cf9d977207/grh-compact/src/New/Rows/RowVolume.lean\#L67-L70}{{\leanfont Data.\allowbreak{}groupAlgAction\_auxLift}}}
\expandafter\newcommand\csname LN@rangeGoodLift\endcsname{\href{https://github.com/lsernau/grh-lean/blob/72fead102f663d5c7a407964a13dc0cf9d977207/grh-compact/src/New/Rows/RowVolume.lean\#L75-L82}{{\leanfont Data.\allowbreak{}range\_goodLift\_le}}}
\newcommand{\ConcordanceRows}{\crow{normDetSq}{\href{https://github.com/lsernau/grh-lean/blob/72fead102f663d5c7a407964a13dc0cf9d977207/grh-compact/src/New/Rows/RowVolume.lean\#L15-L19}{{\leanfont normDet\_sq\_eq\_pow\_mul\_det}}}{RowVolume.lean}{15--19}
\crow{liftVolume}{\href{https://github.com/lsernau/grh-lean/blob/72fead102f663d5c7a407964a13dc0cf9d977207/grh-compact/src/New/Rows/RowVolume.lean\#L23-L23}{{\leanfont Data.\allowbreak{}liftVolume}}}{RowVolume.lean}{23--23}
\crow{chainTheta}{\href{https://github.com/lsernau/grh-lean/blob/72fead102f663d5c7a407964a13dc0cf9d977207/grh-compact/src/New/Rows/RowVolume.lean\#L27-L27}{{\leanfont Data.\allowbreak{}Θ}}}{RowVolume.lean}{27--27}
\crow{chainVolume}{\href{https://github.com/lsernau/grh-lean/blob/72fead102f663d5c7a407964a13dc0cf9d977207/grh-compact/src/New/Rows/RowVolume.lean\#L30-L30}{{\leanfont Data.\allowbreak{}𝔷}}}{RowVolume.lean}{30--30}
\crow{volumeProduct}{\href{https://github.com/lsernau/grh-lean/blob/72fead102f663d5c7a407964a13dc0cf9d977207/grh-compact/src/New/Rows/RowVolume.lean\#L37-L58}{{\leanfont Data.\allowbreak{}volume\_iota\_eq\_prod\_liftVolume}}}{RowVolume.lean}{37--58}
\crow{rangeTupleIdem}{\href{https://github.com/lsernau/grh-lean/blob/72fead102f663d5c7a407964a13dc0cf9d977207/grh-compact/src/New/Rows/RowVolume.lean\#L62-L65}{{\leanfont Data.\allowbreak{}range\_tupleIdem\_le}}}{RowVolume.lean}{62--65}
\crow{groupAlgActionAuxLift}{\href{https://github.com/lsernau/grh-lean/blob/72fead102f663d5c7a407964a13dc0cf9d977207/grh-compact/src/New/Rows/RowVolume.lean\#L67-L70}{{\leanfont Data.\allowbreak{}groupAlgAction\_auxLift}}}{RowVolume.lean}{67--70}
\crow{rangeGoodLift}{\href{https://github.com/lsernau/grh-lean/blob/72fead102f663d5c7a407964a13dc0cf9d977207/grh-compact/src/New/Rows/RowVolume.lean\#L75-L82}{{\leanfont Data.\allowbreak{}range\_goodLift\_le}}}{RowVolume.lean}{75--82}}
% the notations of the defined functions: one slot for each explicit argument of the Lean declaration
\newcommand{\liftVolume}[3]{\operatorname{vol}_{#1}(#2,#3)}
\newcommand{\chainTheta}[2]{\Theta_{#1,#2}}
\newcommand{\chainVolume}[2]{\mathfrak{z}_{#1}(#2)}
\newcommand{\UnitMembers}{219}
\newcommand{\UnitInNote}{8}
\newcommand{\UnitProofs}{85}
\newcommand{\UnitSteps}{634}
\newcommand{\UnitAllSteps}{680}


\linespread{1.08}                       % 11 pt on 14.7 pt
\setlength{\parindent}{0pt}
\setlength{\parskip}{5.5pt plus 1.5pt minus 0.5pt}
\setlength{\abovedisplayskip}{7.5pt plus 2pt minus 3pt}
\setlength{\belowdisplayskip}{7.5pt plus 2pt minus 3pt}
\setlength{\abovedisplayshortskip}{3pt plus 2pt}
\setlength{\belowdisplayshortskip}{6pt plus 2pt minus 2pt}
\emergencystretch=1.5em
\clubpenalty=10000 \widowpenalty=10000 \displaywidowpenalty=10000
\relpenalty=7000 \binoppenalty=9000
\raggedbottom
\pagestyle{fancy}\fancyhf{}\renewcommand{\headrulewidth}{0pt}
\fancyfoot[C]{\sffamily\fontsize{7.6}{9}\selectfont\color{inkfaint}\thepage\,/\,\pageref*{LastPage}}
\setlist{topsep=2pt,partopsep=0pt,parsep=0pt,itemsep=4pt,leftmargin=2.1em,labelsep=0.55em}

\newcommand{\Q}{\mathbb{Q}}
\newcommand{\Z}{\mathbb{Z}}
\newcommand{\F}{\mathbb{F}}
\newcommand{\len}{\operatorname{len}}
\newcommand{\cok}{\operatorname{cok}}
\newcommand{\End}{\operatorname{End}}
\newcommand{\Hom}{\operatorname{Hom}}
\newcommand{\Gal}{\operatorname{Gal}}
\newcommand{\diag}{\operatorname{diag}}
\newcommand{\kw}[1]{\emph{#1}}
\newcommand{\m}[1]{\mbox{$#1$}}        % a short formula that must not be broken across lines
\newcommand{\head}[1]{\par\addvspace{17pt plus 4pt minus 2pt}\Needspace{6\baselineskip}%
  \noindent{\fontsize{13.2}{16}\selectfont\bfseries\color{ink}#1}\par\nopagebreak\vspace{-1pt}}
\newcommand{\sect}[2]{\refstepcounter{section}\label{#2}\par\addvspace{19pt plus 4pt minus 2pt}\Needspace{7\baselineskip}%
  \noindent\llap{{\fontsize{13.2}{16}\selectfont\bfseries\color{inkfaint}\thesection}\hspace{9pt}}%
  {\fontsize{13.2}{16}\selectfont\bfseries\color{ink}#1}\par\nopagebreak\vspace{-1pt}}
\newcommand{\subhead}[1]{\par\addvspace{9pt plus 2pt}\Needspace{4\baselineskip}\noindent\textbf{#1}\enspace\ignorespaces}
% where a statement lives in the Lean archive
\newcommand{\leantag}{{\sffamily\bfseries\fontsize{6.3}{8}\selectfont\color{leancol}\addfontfeatures{LetterSpace=7}LEAN}}
\newcommand{\sep}{\unskip\nobreak\hspace{5pt}{\color{inkfaint}\textperiodcentered}\hspace{5pt plus 2pt}\ignorespaces}
\newcommand{\leanline}[1]{\par\nopagebreak\vspace{-2.5pt}{\fontsize{8.9}{12.4}\selectfont\color{inksoft}\raggedright
  \noindent\llap{\leantag\hspace{9pt}}#1\par}\vspace{1pt}}
% a displayed statement in a tinted box: #1 colour, #2 text
\newcommand{\statement}[2]{\par\addvspace{7pt}\begingroup\setlength{\fboxsep}{8.5pt}\setlength{\fboxrule}{0.6pt}%
  \noindent\fcolorbox{#1!30!white}{#1!3!white}{\parbox{\dimexpr\linewidth-2\fboxsep-2\fboxrule\relax}{#2}}%
  \endgroup\par\addvspace{7pt}}
\newcommand{\nt}[2]{\par\noindent\hangindent=1.1em\hangafter=1 #1\enspace{\color{inksoft}#2}\par\vspace{1pt}}
\newcommand{\legdot}[1]{\tikz[baseline=-0.62ex]\fill[#1] (0,0) circle (2.5pt);\hspace{3pt}}
\newcommand{\crow}[4]{\csname CS@#1\endcsname & #2 & \textcolor{inksoft}{#3} & \textcolor{inksoft}{#4}\\}
\newcommand{\leanquoteline}[2]{\par\noindent\hangindent=\dimexpr#1\fontdimen2\font+2.2em\relax\hangafter=1
  \hspace*{#1\fontdimen2\font}#2}
\newcommand{\quoted}[1]{\par\vspace{3pt}{\ttfamily\fontsize{7.5}{10.4}\selectfont\color{leancol}\raggedright
  \setlength{\parskip}{0pt}\leanquote{#1}\par}\vspace{4pt}}
\newcommand{\bx}[3]{{\sffamily\bfseries\fontsize{7.3}{9.2}\selectfont\color{#1!82!black}#2}\\[1.5pt]#3}
\newcommand{\cs}[2]{\expandafter\newcommand\csname CS@#1\endcsname{#2}}

% theorems
\newtheoremstyle{plainx}{7pt plus 2pt}{7pt plus 2pt}{\itshape}{}{\bfseries}{.}{0.5em}{}
\newtheoremstyle{uprightx}{7pt plus 2pt}{7pt plus 2pt}{}{}{\bfseries}{.}{0.5em}{}
\theoremstyle{plainx}
\newtheorem{theorem}{Theorem}[section]
\newtheorem{lemma}[theorem]{Lemma}
\newtheorem{proposition}[theorem]{Proposition}
\newtheorem{corollary}[theorem]{Corollary}
\theoremstyle{uprightx}
\newtheorem{definition}[theorem]{Definition}
\newtheorem{setting}[theorem]{Setting}          % hypotheses shared by several statements; each of them cites it by number
\newtheorem{remark}[theorem]{Remark}
\newtheorem{example}[theorem]{Example}
\renewcommand{\qedsymbol}{{\setlength{\fboxsep}{0pt}\setlength{\fboxrule}{0.5pt}\fbox{\rule{0pt}{0.46em}\hspace{0.46em}}}}
% what holds throughout the note: the ambient objects only (implicit and instance arguments), never a hypothesis
\newenvironment{conventions}{\par\addvspace{5pt}\noindent\textbf{Conventions.}\enspace\ignorespaces}{\par\addvspace{3pt}}
\makeatletter
\renewenvironment{proof}[1][Proof]{\par\pushQED{\qed}\normalfont\topsep0pt\partopsep0pt
  \trivlist\item[\hskip\labelsep\itshape #1.]\ignorespaces}{\popQED\endtrivlist\@endpefalse}
\makeatother

% where each declaration of the concordance is treated: \cs{key}{place}, the place naming a \label or a lettered item
\cs{normDetSq}{Lemma~\ref{l:gram}}
\cs{liftVolume}{Definition~\ref{d:vol}}
\cs{chainTheta}{Definition~\ref{d:vol}}
\cs{chainVolume}{Definition~\ref{d:vol}}
\cs{volumeProduct}{Theorem A}
\cs{rangeTupleIdem}{Lemma~\ref{l:idem}}
\cs{groupAlgActionAuxLift}{Lemma~\ref{l:equiv}}
\cs{rangeGoodLift}{Lemma~\ref{l:range}}
\newcommand{\CC}{\mathbb{C}}
\newcommand{\ZZ}{\mathbb{Z}}
\newcommand{\RR}{\mathbb{R}}
\newcommand{\AutC}{\operatorname{Aut}(\CC)}
\newcommand{\Row}{\mathrm{Row}}
\newcommand{\Mf}{\mathcal{M}}
\newcommand{\Mg}{\mathcal{M}^{\circ}}
\newcommand{\Ff}{\mathcal{F}}
\newcommand{\alift}{\operatorname{lift}^{*}}
\newcommand{\chain}{\operatorname{chain}}
\newcommand{\Hk}{\mathbb{T}}
\newcommand{\sfG}{\mathsf{G}}
\newcommand{\idem}{\operatorname{idem}}
\newcommand{\vol}{\operatorname{vf}}
\newcommand{\ip}[2]{\langle #1,#2\rangle}
\newcommand{\inputbox}[1]{\statement{inksoft}{#1}}

\begin{document}
\color{ink}

{\raggedright
{\sffamily\bfseries\fontsize{7.2}{9}\selectfont\color{leancol}\addfontfeatures{LetterSpace=8}READBACK NOTE · CHAPTER VII · UNIT 143 · PART 1 OF 4\par}
\vspace{5pt}
{\fontsize{21.5}{25.5}\selectfont\bfseries The volume of the chain of theta lifts as a product of lift volumes\par}
\vspace{5pt}
{\fontsize{12.4}{16}\selectfont\itshape\color{inksoft} One unit of the Lean artifact \ArtifactName, read back into
mathematics\par}}
\vspace{8pt}
{\color{rulegrey}\hrule height 0.6pt}
\vspace{7pt}
\thispagestyle{empty}

\textbf{Summary.} The rows of the proof are the good automorphic forms on an indefinite unitary group of rank~$r$,
and the chain of lifts sends them, through $m$ theta lifts, to algebraic modular forms on a definite unitary group.
This note proves that the volume of the chain on the rows, the determinant of its Gram operator, is a constant to
the power of the dimension of the rows times the product of the volumes of $m$ single lifts, each taken at a
definite group of fixed rank and at the matching power of the twist. The constant depends on the twist, but only
through quantities whose ratio to the untwisted one has controlled height. The proof compares Gram forms: by the
artifact's inner product formula, a coset-sum form of the Rallis inner product formula, each Gram form is a constant
times a Hecke operator, and the same Hecke operators occur on both sides.

\statement{leancol}{\fontsize{9.6}{13.2}\selectfont\textbf{What this note is.} The artifact \ArtifactName{} is a
body of Lean code: the compact edition of a formal proof, due to L.~Alpöge with the AI system Claude, of the Riemann
hypothesis and of its generalisation to Dirichlet $L$-functions. That headline is a claim. The authors report that
every file compiles with Lean v4.35.0 against a pinned version of Mathlib and that Lean's kernel accepted the
proofs; an independent build is recorded in the handout's record of checks when it reports. No mathematician has
read the argument through, and nothing here says that the Riemann hypothesis is proved. This note is the first of
four on one unit of the artifact, the folder of the rows; it treats \UnitInNote{} of the unit's \UnitMembers{}
declarations, those of the file on the volume of the chain. Theorem~A and Lemma~\ref{l:gram} are proved in full from
Inputs 1 to~7, which belong to other units and are stated in a box and not proved here; Lemmas \ref{l:idem}
to~\ref{l:range} are proved in full from two further facts of Chapter~VIII, which their proofs name. Every proof of
this note is complete relative to the inputs it names. Every
statement below was compared with its Lean text, argument by argument. The Lean text read is \PinText; Appendix~A
is the concordance of \ConcDecls{} declarations.}

\sect{The statement}{s:intro}
Theorem~C of the handout bounds a height from below by comparing the chain of lifts at the trivial twist with the
chain at well-chosen twists. It needs to know how the volume of the chain depends on the twist. The answer has two
parts: a factorisation, which is Theorem~A here, and a bound for each factor, which is the subject of the note
\emph{The height of a ratio of lift volumes}. The objects are those of Setting~\ref{set:data} and
Definition~\ref{d:vol}.

\statement{arith}{\phantomsection\label{t:A}\textbf{Theorem A} (the volume of the chain is a product).\enspace
\emph{In Setting~\ref{set:data}, for all parameters~$P$ and all admissible data~$D$ for~$P$, with number of
lifts~$m$, there is a function~$c$ from the $m$-fold twists of~$D$ to~$\CC$ with the following two properties.
\begin{enumerate}[label=\textup{(\alph*)},leftmargin=2.3em,itemsep=2pt]
\item For every family $\theta=(\theta_{\ell,\iota})$ of $m$-fold twists of~$D$ that is fit for~$D$, the family
\m{(\ell,\iota)\mapsto|c(\theta_{\ell,\iota})/c(1)|^{-2}} has height at most \m{m^2\cdot2\,t(D)\,x^{999}}, where $x$ is
the first scale of~$D$ and \m{t(D)=5r}.
\item For every $m$-fold twist~$\vartheta$ of~$D$,
\[\chainVolume{D}{\vartheta}\ne0\qquad\text{and}\qquad
\chainVolume{D}{\vartheta}=|c(\vartheta)|^{\dim\Row(D)}\cdot\prod_{j=1}^{m}\liftVolume{D}{\vartheta}{j}.\]
\end{enumerate}}}
\leanline{\lean{volumeProduct}\sep property (a) is the formal text's ``$c$ costs \m{m\cdot m} squares''
(Input~7)}

Here $\chainVolume{D}{\vartheta}$ is the volume of the chain $\chainTheta{D}{\vartheta}$ of $m$ lifts on the rows,
and $\liftVolume{D}{\vartheta}{j}$ is the volume of the $j$-th single lift, formed with the twist~$\vartheta^{j}$,
on the good forms of its own source. The exponent \m{\dim\Row(D)} appears because all the volumes are determinants
on spaces of this one dimension. Since \m{\dim\Row(D)>0} (Theorem~B of the handout), property~(b) gives
\m{c(\vartheta)\ne0} for every $m$-fold twist~$\vartheta$, so that the family in~(a) is defined. The formal text does
not need this remark, since it reads a quotient by zero as zero.

\sect{The setting}{s:setting}
\begin{conventions}\phantomsection\label{c:conventions}
Inner product spaces are complex, and an inner product \m{\ip{y'}{y}} is conjugate-linear in~$y'$ and linear
in~$y$.
For a linear map~$u$ between finite-dimensional inner product spaces, $u^{*}$ is its adjoint. The implicit
arguments of the declarations of this note are the parameters~$P$ (in all of them) and, in
Lemmas~\ref{l:gram} and~\ref{l:idem}, an inner product space~$U$ and, in Lemma~\ref{l:gram}, a second one~$U'$
and the scalar~$a$; in Lemma~\ref{l:idem} also the homomorphism~$\varrho$ and the tuple~$\mathbf{c}$; in
Lemmas \ref{l:equiv} and~\ref{l:range} the tame character~$\vartheta$ and the lift index~$j$, in the first of them the
element~$\xi$ and in the second the subspace~$\mathcal{N}$. Each of them is introduced again in the statement that
uses it.
\end{conventions}

\begin{setting}[data, forms and lifts; all taken from other units]\label{set:data}
\emph{Data.} $P=(N_0,d)$ are parameters and $D$ is admissible data for~$P$ (Setting~1 of the overview of the
handout; the notes on Chapter~III). Attached to~$D$ are the CM field \m{E(D)} with its ring of
integers~$\mathcal{O}$, the number~$m$ of lifts, the rank~$r$, the step \m{t(D)=5r}, the first scale~$x$, the second
level \m{g(D)=g_0g_1} and the level \m{N(D)}, a prime (Setting~2 of the overview), the finite set of \kw{chosen
primes} of \m{E(D)}, and the finite group $\sfG(D)$, the product of the groups $\mathrm{GL}_r$ of the residue fields of the chosen primes. A
\kw{lift index} is an integer~$j$ with \m{1\le j\le m}.

\emph{Spaces.} $Y_0(D),\dots,Y_m(D)$ are hermitian spaces over \m{E(D)}, the space $Y_i(D)$ of rank \m{r+i\,t(D)};
$Y_0(D)$ is indefinite and $Y_m(D)$ is definite. For a lift index~$j$, $W_j(D)$ is a diagonal hermitian space of
rank~$r$.

\emph{Forms of principal level.} For a hermitian space~$V$ of rank~$r$ over \m{E(D)}, $\Mf_D(V)$ is the space of
automorphic forms on the unitary group of~$V$ of a fixed type at the infinite places and fixed by the principal
congruence subgroup at the chosen primes. It is an inner product space under the Petersson product, the group
algebra \m{\CC[\sfG(D)]} acts on it, and it is finite-dimensional for \m{V=Y_0(D)} and for \m{V=W_j(D)}.

\emph{Good forms and rows.} A \kw{tuple of labels} is a family of monic polynomials of degree~$r$, one over the
residue field of each chosen prime. Each tuple~$\mathbf{c}$ has a central idempotent \m{\idem(\mathbf{c})} of
\m{\CC[\sfG(D)]}; idempotents of different tuples are orthogonal. A finite set of tuples is called \kw{good}, and
$e_0(D)$ is the sum of \m{\idem(\mathbf{c})} over the good tuples~$\mathbf{c}$; it is a central idempotent. The
\kw{good forms} on~$V$ are \m{\Mg_D(V)=e_0(D)\,\Mf_D(V)}, and the \kw{rows} are \m{\Row(D)=\Mg_D(Y_0(D))}.

\emph{Hecke operators.} $\Hk_D$ is the Hecke algebra of~$D$: an algebra whose elements~$T$ act on every
$\Mf_D(V)$, by an endomorphism~$T_V$ that preserves $\Mg_D(V)$; it contains the Hecke operators at the primes
away from the level. A \kw{Hecke eigensystem} is a function~$p$ that assigns a complex number to each of these
Hecke operators, and \m{\Mf_D(V)[p]} is the space of the forms on which each of them acts by the value of~$p$.

\emph{Twists.} A \kw{tame character} of~$D$ is a character~$\vartheta$ of the norm-one ideles of \m{E(D)}, trivial
on the norm-one elements of \m{E(D)} and at the infinite places, at most tamely ramified, with a finite
set~\m{Q(\vartheta)} of primes outside of which it is unramified; \m{Q(\vartheta)} contains the complex conjugate
of none of its members, and the norms of its members are prime to \m{g(D)N(D)}. It is an \kw{$m$-fold twist} if
$\vartheta^{n}$ is ramified at every prime of \m{Q(\vartheta)} for \m{1\le n\le m}. The trivial character~$1$ is
an $m$-fold twist with \m{Q(1)=\emptyset}, and every power $\vartheta^{n}$ is a tame character with
\m{Q(\vartheta^{n})=Q(\vartheta)}. For an integer \m{e\ge0}, \m{\mathfrak{m}_D(\vartheta,e)} is the ideal of
\m{E(D)} equal to the $e$-th power of the product of the chosen primes times the product of the primes in
\m{Q(\vartheta)}, and \m{|\mathcal{O}/\mathfrak{m}_D(\vartheta,e)|} is the number of elements of its residue ring.

\emph{Single lifts.} For a tame character~$\vartheta$ and a lift index~$j$, the \kw{auxiliary pair}
\m{O_D(\vartheta,j)} consists of $W_j(D)$, of a definite hermitian space of rank \m{r+t(D)} and of a character
built from~$\vartheta$; \m{\Ff(O_D(\vartheta,j))} is its space of algebraic modular forms on the definite space, a
finite-dimensional inner product space on which \m{\CC[\sfG(D)]} acts (the note \emph{Auxiliary dual pairs,
determinant polynomials and the carrier}). The \kw{adjoint lift}
\[\alift_D(\vartheta,j)\colon\Mf_D(W_j(D))\longrightarrow\Ff(O_D(\vartheta,j))\]
is the adjoint of a theta lift from the definite space to~$W_j(D)$.

\emph{The chain.} For an $m$-fold twist~$\vartheta$, the linear map
\m{\chain_D(\vartheta)\colon\Mf_D(Y_0(D))\to\mathcal{A}_{D,\vartheta}} is the composite of $m$ theta lifts
\m{Y_0(D)\to Y_1(D)\to\dots\to Y_m(D)}, the $i$-th formed with~$\vartheta^{i}$, with values in a finite-dimensional
inner product space $\mathcal{A}_{D,\vartheta}$ of algebraic modular forms on the definite space~$Y_m(D)$.

\emph{Heights and fit families.} A family \m{Z=(Z_{\ell,\iota})} of complex numbers, indexed by the pairs of a
prime~$\ell$ and a field isomorphism~$\iota$ from~$\CC$ to an algebraic closure of~$\mathbb{Q}_\ell$, has
\kw{height at most~$h$}, for a real number~$h$, if there are \m{H\in\CC} with \m{H\ne0}, a real number~$\alpha$
and real numbers~$\eta_\ell$ with \m{|\sigma(H)|\le e^{\alpha}} for all \m{\sigma\in\AutC},
\m{|\iota(H\,Z_{\ell,\iota})|_\ell\le e^{\eta_\ell}} for all \m{(\ell,\iota)}, and such that $\alpha$ plus the sum of
the~$\eta_\ell$ over the primes~$\ell$ below any sufficiently large bound is at most~$h$. A family
\m{\theta=(\theta_{\ell,\iota})} of $m$-fold twists of~$D$, with the same index set, may be \kw{fit for~$D$}; the
definition belongs to Chapter~VIII and enters this note only through Input~7.
\end{setting}

\begin{definition}[volumes]\label{d:vol}
For a linear map~$u$ between finite-dimensional inner product spaces, the \kw{volume factor} \m{\vol(u)} is the
real number \m{\ge0} with \m{\vol(u)^2=\det(u^{*}u)}; it vanishes exactly when $u$ is not injective. In
Setting~\ref{set:data}:
\begin{enumerate}[label=\textup{(\roman*)},leftmargin=2.3em,itemsep=2pt]
\item for an $m$-fold twist~$\vartheta$ of~$D$ and a lift index~$j$, the \kw{lift volume} is
\[\liftVolume{D}{\vartheta}{j}=\vol(u)^2=\det(u^{*}u),\]
where \m{u\colon\Mg_D(W_j(D))\to\Ff(O_D(\vartheta^{j},j))} is the restriction of \m{\alift_D(\vartheta^{j},j)}, the
adjoint lift for the twist~$\vartheta^{j}$, to the good forms;
\item for an $m$-fold twist~$\vartheta$ of~$D$, the \kw{chain of lifts on the rows} is the restriction
\[\chainTheta{D}{\vartheta}=\chain_D(\vartheta)\big|_{\Row(D)}\colon\ \Row(D)\to\mathcal{A}_{D,\vartheta};\]
\item for an $m$-fold twist~$\vartheta$ of~$D$, the \kw{volume of the chain} is
\[\chainVolume{D}{\vartheta}=\vol(\chainTheta{D}{\vartheta})^2=\det(\chainTheta{D}{\vartheta}^{*}\,\chainTheta{D}{\vartheta}).\]
\end{enumerate}
Both volumes are real numbers \m{\ge0}, regarded as complex numbers.
\end{definition}
\leanline{\lean{liftVolume}\sep \lean{chainTheta}\sep \lean{chainVolume}}

\inputbox{\textbf{Inputs from other units.} The seven results below are taken for granted. They are stated as the
dossier of this unit reads them; the notes on their units prove them. In all of them $P$ and~$D$ are as in
Setting~\ref{set:data}.

\vspace{4pt}\textbf{Input 1} (unit \emph{Exchange}; the inner product formula for one lift).\enspace For every lift
index~$j$ there is \m{A_j\in\CC}, \m{A_j\ne0}, such that for every tame character~$\vartheta$ of~$D$ the map
\m{\alift_D(\vartheta,j)} is injective and, for all \m{y',y\in\Mf_D(W_j(D))},
\[\ip{\alift_D(\vartheta,j)\,y'}{\alift_D(\vartheta,j)\,y}=A_j\,\delta_D(\vartheta)\sum_{\gamma}R_{D}(\vartheta,j,W_j(D);y',y;\gamma).\]
Here the terms \m{R_{D}(\vartheta,j,V;y',y;\gamma)}, defined for every hermitian space~$V$ of rank~$r$ and all
\m{y',y\in\Mf_D(V)}, are those of the artifact's coset-sum form of the Rallis inner product formula (a matrix
coefficient of a translation at a coset~$\gamma$, weighted by a local factor that depends on $\vartheta$ and~$j$;
this note calls them Rallis terms), the series converges
unconditionally, and \m{\delta_D(\vartheta)\ne0} is the \kw{Gram factor}: a ratio of indices and orders of finite
groups that depends only on the chosen primes and on \m{Q(\vartheta)}.

\vspace{3pt}\textbf{Input 2} (unit \emph{Exchange}; the first lift of each row of the triangle).\enspace For every
$m$-fold twist~$\vartheta$ there are a family of inner product spaces of forms \m{F_\vartheta(i,n)}, for
\m{i+n\le m}, with linear maps \m{a_\vartheta(i,n)\colon F_\vartheta(i,n)\to F_\vartheta(i,n+1)}, which are theta
lifts, isometric embeddings \m{\jmath_i\colon\Mf_D(Y_0(D))\to F_\vartheta(i,0)}, and positive integers
\m{N_n(\vartheta)}, for \m{0\le n\le m}. For every \m{i<m} there is \m{C_i\in\CC} such that for every $m$-fold
twist~$\vartheta$ and all \m{z,y\in\Mf_D(Y_0(D))}
\[N_0(\vartheta)\,\ip{a_\vartheta(i,0)\jmath_i z}{a_\vartheta(i,0)\jmath_i y}=C_i\sum_{\gamma}R_{D}(\vartheta^{i+1},i+1,Y_0(D);z,y;\gamma).\]

\vspace{3pt}\textbf{Input 3} (unit \emph{Exchange}; the Gram form of the weighted chain).\enspace Let
\m{u_\vartheta\colon F_\vartheta(0,0)\to F_\vartheta(0,m)} be the composite of the maps
\m{N_n(\vartheta)\,a_\vartheta(0,n)}, \m{0\le n<m}. Suppose that for every $m$-fold twist~$\vartheta$ and every
\m{i<m} an endomorphism \m{S(\vartheta,i)} of \m{\Row(D)} is given with
\m{N_0(\vartheta)\ip{a_\vartheta(i,0)\jmath_i z}{a_\vartheta(i,0)\jmath_i y}=\ip{z}{S(\vartheta,i)y}} for all
\m{z\in\Mf_D(Y_0(D))} and \m{y\in\Row(D)}. Then there is a function~$\kappa$ from the $m$-fold twists to~$\CC$ that
costs \m{m\cdot m} squares (Input~7), and for every $m$-fold twist~$\vartheta$ there is an endomorphism \m{T(\vartheta)}
of \m{\Row(D)} with \m{\det T(\vartheta)=\prod_{i<m}\det S(\vartheta,i)} and
\[\ip{u_\vartheta\jmath_0z}{u_\vartheta\jmath_0y}=\kappa(\vartheta)\,\frac{N_0(\vartheta)\,N_1(\vartheta)^{m}}{N_m(\vartheta)}\,\ip{z}{T(\vartheta)y}\]
for all \m{z\in\Mf_D(Y_0(D))} and \m{y\in\Row(D)}.}

\inputbox{\textbf{Inputs from other units, continued.}

\vspace{4pt}\textbf{Input 4} (unit \emph{Exchange}; the chain at its level).\enspace For every $m$-fold
twist~$\vartheta$ the map \m{\chain_D(\vartheta)} is injective. Let \m{I(\vartheta)} be the index of the level of
\m{\mathcal{A}_{D,\vartheta}} in the stabiliser of the standard lattice of~$Y_m(D)$, and let $\omega$ be the dimension
of the weight of \m{\mathcal{A}_{D,\vartheta}}, a positive integer that does not depend on~$\vartheta$. Then
\m{\ip{\chain_D(\vartheta)z}{\chain_D(\vartheta)y}=I(\vartheta)\,\omega\,\ip{u_\vartheta\jmath_0z}{u_\vartheta\jmath_0y}}
for all \m{z,y\in\Mf_D(Y_0(D))}. Moreover there are a number~$G$ and an integer~$l$ with
\[I(\vartheta)=G\cdot|\mathcal{O}/\mathfrak{m}_D(\vartheta,0)|^{\,l}\cdot\delta_D(\vartheta)^{m}\cdot\frac{N_m(\vartheta)}{N_0(\vartheta)\,N_1(\vartheta)^{m}}
\qquad\text{for every $m$-fold twist }\vartheta.\]

\vspace{3pt}\textbf{Input 5} (unit \emph{Trace}; one Hecke element for two Rallis series).\enspace Let
$\vartheta$ be a tame character of~$D$ and $j$ a lift index, let $V$ and~$V'$ be hermitian spaces of rank~$r$ over
\m{E(D)} whose signatures satisfy the gap condition of that unit (the proof below takes \m{V=Y_0(D)} and
\m{V'=W_{i+1}(D)}), let \m{A,A'\in\CC}, and let $B$ and~$B'$ be complex functions of two forms. If
\m{A\sum_\gamma R_D(\vartheta,j,V;z,y;\gamma)=B(z,y)} for all \m{z,y\in\Mf_D(V)} and
\m{A'\sum_\gamma R_D(\vartheta,j,V';z,y;\gamma)=B'(z,y)} for all \m{z,y\in\Mf_D(V')}, the series converging
unconditionally, then there is \m{T\in\Hk_D} with \m{B(z,y)=A\ip{z}{T_Vy}} and \m{B'(z,y)=A'\ip{z}{T_{V'}y}} for all
such $z$ and~$y$.

\vspace{3pt}\textbf{Input 6} (unit \emph{Trace}; good forms on different spaces).\enspace For every lift
index~$j$ and every \m{T\in\Hk_D}, the determinants of~$T$ on \m{\Mg_D(Y_0(D))} and on \m{\Mg_D(W_j(D))} are equal,
and \m{\dim\Row(D)=\dim\Mg_D(W_j(D))}.

\vspace{3pt}\textbf{Input 7} (unit \emph{Exchange}; the cost of squares).\enspace For an integer \m{n\ge0}, a
function~$c$ from the $m$-fold twists of~$D$ to~$\CC$ \kw{costs $n$ squares} if, for every family~$\theta$ of
$m$-fold twists that is fit for~$D$, the family \m{(\ell,\iota)\mapsto|c(\theta_{\ell,\iota})/c(1)|^{-2}} has height
at most \m{n\cdot2\,t(D)\,x^{999}}. (i) If $c$ costs $n$ squares, $c'$ costs $n'$ squares and \m{n+n'\le n''}, then
the product \m{c\,c'} costs $n''$ squares. (ii) If there are a real number~$G'$, an integer \m{e\ge0} and an
integer~$l'$ with \m{|c(\vartheta)|^2=G'\cdot|\mathcal{O}/\mathfrak{m}_D(\vartheta,e)|^{\,l'}} for every $m$-fold
twist~$\vartheta$, then $c$ costs $0$ squares.}

Property~(a) of Theorem~A says that $c$ costs \m{m\cdot m} squares, and the function~$\kappa$ of Input~3 costs
\m{m\cdot m} squares. Inputs 1 and~2 rest on the artifact's coset-sum (doubling) form of the Rallis inner product
formula for a theta lift~$L$ in the stable range,
\[\ip{Ly'}{Ly}=A\,c_F\sum_\gamma\rho(\gamma)\,\ip{y'}{R(\gamma)y},\]
with a constant \m{A\ne0}, the mass~$c_F$ of the finite datum, a right translation \m{R(\gamma)} and a matrix
coefficient \m{\rho(\gamma)} of the vacuum. For it the catalogue of the
handout cites S.~Rallis (\emph{Injectivity properties of liftings associated to Weil representations}, Compositio
Math.\ 52 (1984), §1 and Theorem~1.6) and J.-S.~Li (\emph{Non-vanishing theorems for the cohomology of certain
arithmetic quotients}, J.\ reine angew.\ Math.\ 428 (1992), §2, Theorem~2.1, eq.~(26), p.~184). The second reading
classes these as analogues: that the formal statement is a special case of the cited ones is not supported by the
citations, and the uniform constant, the measures, the evaluation as a coset sum and the injectivity would each
need an argument. So the note does not call Inputs 1 and~2 instances of the classical formula. In Input~2 the spaces \m{F_\vartheta(i,n)} form a triangle whose row \m{i=0} is the chain
\m{Y_0(D)\to\dots\to Y_m(D)}; the other rows serve the proof of Input~3 only.

\sect{Proofs}{s:proofs}
\begin{lemma}[a Gram form determines the volume]\label{l:gram}
Let $U$ and~$U'$ be finite-dimensional inner product spaces, \m{u\colon U\to U'} a linear map, $S$ an endomorphism
of~$U$ and \m{a\in\CC}. If \m{\ip{uy'}{uy}=a\,\ip{y'}{Sy}} for all \m{y',y\in U}, then
\[\vol(u)^2=a^{\dim U}\cdot\det S.\]
\end{lemma}
\begin{proof}
For all \m{y',y\in U},
\[\ip{y'}{u^{*}uy}=\ip{uy'}{uy}=a\,\ip{y'}{Sy}=\ip{y'}{(aS)y}.\]
So \m{u^{*}u=aS}, and
\m{\vol(u)^2=\det(u^{*}u)=\det(aS)=a^{\dim U}\det S}.
\end{proof}
\leanline{\lean{normDetSq}}

\begin{proof}[Proof of Theorem A]
Fix $P$ and~$D$, and write \m{n_0=\dim\Row(D)}.

\emph{The single lifts.} Take the numbers \m{A_j\ne0} of Input~1 and the numbers~$C_i$ of Input~2. Let
$\vartheta$ be an $m$-fold twist and \m{0\le i<m}. The power \m{\vartheta^{i+1}} is a tame character with
\m{Q(\vartheta^{i+1})=Q(\vartheta)}, so \m{\delta_D(\vartheta^{i+1})=\delta_D(\vartheta)}. Input~1, for the tame
character \m{\vartheta^{i+1}} and the lift index \m{i+1}, and Input~2 give two series of Rallis terms with the same
parameters \m{(\vartheta^{i+1},i+1)}, one on \m{W_{i+1}(D)} with the constant \m{A_{i+1}\delta_D(\vartheta)} and one
on $Y_0(D)$ with the constant~$C_i$. By Input~5 there is \m{T_{\vartheta,i}\in\Hk_D} with
\begin{align*}
N_0(\vartheta)\,\ip{a_\vartheta(i,0)\jmath_iz}{a_\vartheta(i,0)\jmath_iy}&=C_i\,\ip{z}{T_{\vartheta,i}\,y},\\
\ip{\alift_D(\vartheta^{i+1},i+1)\,z}{\alift_D(\vartheta^{i+1},i+1)\,y}&=A_{i+1}\,\delta_D(\vartheta)\,\ip{z}{T_{\vartheta,i}\,y},
\end{align*}
the first for all \m{z,y\in\Mf_D(Y_0(D))} and the second for all \m{z,y\in\Mf_D(W_{i+1}(D))}, where $T_{\vartheta,i}$
acts on the forms of the space in question. Since $T_{\vartheta,i}$ preserves the good
forms, the second identity and Lemma~\ref{l:gram}, applied to the restriction of the adjoint lift to
\m{\Mg_D(W_{i+1}(D))}, a space of dimension~$n_0$ by Input~6, give
\begin{equation}\label{e:single}
\liftVolume{D}{\vartheta}{i+1}=\bigl(A_{i+1}\,\delta_D(\vartheta)\bigr)^{n_0}\cdot\det\bigl(T_{\vartheta,i}\mid\Mg_D(W_{i+1}(D))\bigr).
\end{equation}

\emph{The chain.} Let \m{S(\vartheta,i)} be \m{C_i} times the restriction of~$T_{\vartheta,i}$ to \m{\Row(D)}. By
the first identity the hypothesis of Input~3 holds. Take $\kappa$ and the endomorphisms \m{T(\vartheta)} from
Input~3. With Input~4, for \m{z,y\in\Row(D)},
\begin{align*}
\ip{\chainTheta{D}{\vartheta}z}{\chainTheta{D}{\vartheta}y}
&=I(\vartheta)\,\omega\,\kappa(\vartheta)\,\frac{N_0(\vartheta)N_1(\vartheta)^{m}}{N_m(\vartheta)}\,\ip{z}{T(\vartheta)y}\\
&=G\,|\mathcal{O}/\mathfrak{m}_D(\vartheta,0)|^{\,l}\,\delta_D(\vartheta)^{m}\,\omega\,\kappa(\vartheta)\,\ip{z}{T(\vartheta)y}.
\end{align*}
By Lemma~\ref{l:gram} and the formula for \m{\det T(\vartheta)},
\begin{equation}\label{e:chain}
\chainVolume{D}{\vartheta}=\bigl(G\,|\mathcal{O}/\mathfrak{m}_D(\vartheta,0)|^{\,l}\,\delta_D(\vartheta)^{m}\,\omega\,\kappa(\vartheta)\bigr)^{n_0}
\cdot\prod_{i<m}C_i^{\,n_0}\det\bigl(T_{\vartheta,i}\mid\Row(D)\bigr).
\end{equation}

\emph{Comparison.} By Input~6 the determinants of $T_{\vartheta,i}$ on \m{\Row(D)=\Mg_D(Y_0(D))} and on
\m{\Mg_D(W_{i+1}(D))} agree. Insert this into \eqref{e:chain} and write
\[C_i^{\,n_0}\,\delta_D(\vartheta)^{n_0}=(C_i/A_{i+1})^{n_0}\,(A_{i+1}\,\delta_D(\vartheta))^{n_0},\]
which is allowed since \m{A_{i+1}\ne0}. By the identities \eqref{e:single} for \m{i<m},
\[\chainVolume{D}{\vartheta}=c(\vartheta)^{n_0}\cdot\prod_{j=1}^{m}\liftVolume{D}{\vartheta}{j},\qquad
c(\vartheta)=\kappa(\vartheta)\cdot\prod_{i<m}\frac{C_i}{A_{i+1}}\cdot\omega\cdot G\cdot|\mathcal{O}/\mathfrak{m}_D(\vartheta,0)|^{\,l}.\]
The volumes are real and \m{\ge0}. Taking absolute values on both sides gives the formula of~(b), with
\m{|c(\vartheta)|^{n_0}} in the place of \m{c(\vartheta)^{n_0}}.

\emph{Non-vanishing.} The map \m{\chain_D(\vartheta)} is injective by Input~4, so its restriction
\m{\chainTheta{D}{\vartheta}} is injective and \m{\chainVolume{D}{\vartheta}\ne0}.

\emph{The cost.} Write \m{c=\kappa\cdot c_0}. The function~$c_0$ satisfies
\m{|c_0(\vartheta)|^2=G'\cdot|\mathcal{O}/\mathfrak{m}_D(\vartheta,0)|^{2l}} with a real number~$G'$ that does not
depend on~$\vartheta$, so it costs $0$ squares by Input~7\,(ii). The function~$\kappa$ costs \m{m\cdot m}
squares. By Input~7\,(i), with \m{m\cdot m+0\le m\cdot m}, the function~$c$ costs \m{m\cdot m} squares, which
is~(a).
\end{proof}
\leanline{\lean{volumeProduct}\sep the formal proof takes absolute values of the two identities given by
Lemma~\ref{l:gram} before it compares them}

The last three results of the file describe the image of one adjoint lift on the good forms. They are not used
for Theorem~A; the note \emph{The height of a ratio of lift volumes} and Chapter~VIII use them
(§\ref{s:used}).

\begin{lemma}[a good tuple acts inside the good part]\label{l:idem}
In Setting~\ref{set:data}, let $U$ be an inner product space, let
\m{\varrho\colon\CC[\sfG(D)]\to\operatorname{End}_{\CC}(U)} be a homomorphism of $\CC$-algebras, and let
$\mathbf{c}$ be a tuple of labels that is good. Then the image of \m{\varrho(\idem(\mathbf{c}))} is contained in
the image of \m{\varrho(e_0(D))}.
\end{lemma}
\begin{proof}
The idempotents of the tuples are orthogonal and $e_0(D)$ is their sum over the good tuples, of which
$\mathbf{c}$ is one. So \m{\idem(\mathbf{c})=\idem(\mathbf{c})\,e_0(D)=e_0(D)\,\idem(\mathbf{c})}, the second
equality because $e_0(D)$ is central. Hence
\m{\varrho(\idem(\mathbf{c}))=\varrho(e_0(D))\circ\varrho(\idem(\mathbf{c}))}, and the image of a composite lies in
the image of its outer map.
\end{proof}
\leanline{\lean{rangeTupleIdem}}

\begin{lemma}[the adjoint lift commutes with the finite group]\label{l:equiv}
In Setting~\ref{set:data}, let $\vartheta$ be a tame character of~$D$, $j$ a lift index and
\m{\xi\in\CC[\sfG(D)]}. For every \m{y\in\Mf_D(W_j(D))},
\[\alift_D(\vartheta,j)(\xi\cdot y)=\xi\cdot\alift_D(\vartheta,j)(y),\]
where $\xi$ acts on the left through the action of \m{\CC[\sfG(D)]} on \m{\Mf_D(W_j(D))} and on the right through
its action on \m{\Ff(O_D(\vartheta,j))}.
\end{lemma}
\begin{proof}
Both actions are linear in~$\xi$, and $\xi$ is a finite combination of elements of \m{\sfG(D)}. For a group element
the identity is the equivariance of the adjoint lift, which the unit of Chapter~VIII that defines the adjoint
lift proves.
\end{proof}
\leanline{\lean{groupAlgActionAuxLift}}

\begin{lemma}[the image of the adjoint lift on the good forms]\label{l:range}
In Setting~\ref{set:data}, let $\vartheta$ be a tame character of~$D$, $j$ a lift index and $\mathcal{N}$ a linear
subspace of \m{\Ff(O_D(\vartheta,j))}. Suppose that for every Hecke eigensystem~$p$, every
\m{y\in\Mg_D(W_j(D))\cap\Mf_D(W_j(D))[p]} and every good tuple of labels~$\mathbf{c}$: if
\m{\alift_D(\vartheta,j)(y)} lies in the image of the action of \m{\idem(\mathbf{c})} on
\m{\Ff(O_D(\vartheta,j))}, then \m{\alift_D(\vartheta,j)(y)\in\mathcal{N}}. Then
\[\alift_D(\vartheta,j)\bigl(\Mg_D(W_j(D))\bigr)\subseteq\mathcal{N}.\]
\end{lemma}
\begin{proof}
The space \m{\Mf_D(W_j(D))} is the sum of the spaces \m{\Mf_D(W_j(D))[p]} over the Hecke eigensystems~$p$ (a fact of
the unit of Chapter~VIII that defines them). The good forms are the forms \m{e_0(D)\,y}. So it suffices to show, for every~$p$ and every
\m{y\in\Mf_D(W_j(D))[p]}, that the lift of \m{e_0(D)\,y} lies in~$\mathcal{N}$. Now
\m{e_0(D)\,y=\sum_{\mathbf{c}}\idem(\mathbf{c})\,y}, the sum over the good tuples. Fix one. The form
\m{y_{\mathbf{c}}=\idem(\mathbf{c})\,y} is good by Lemma~\ref{l:idem}, applied to the action on \m{\Mf_D(W_j(D))},
and it lies in \m{\Mf_D(W_j(D))[p]}, because the Hecke operators commute with the action of \m{\sfG(D)}. By
Lemma~\ref{l:equiv}, \m{\alift_D(\vartheta,j)(y_{\mathbf{c}})=\idem(\mathbf{c})\cdot\alift_D(\vartheta,j)(y)} lies in the
image of the action of \m{\idem(\mathbf{c})}. The hypothesis gives \m{\alift_D(\vartheta,j)(y_{\mathbf{c}})\in\mathcal{N}},
and the sum over~$\mathbf{c}$ lies in~$\mathcal{N}$.
\end{proof}
\leanline{\lean{rangeGoodLift}}

\sect{Where the results are used}{s:used}
Theorem~A is one of the two inputs of Theorem~C of the handout that the head of the argument names; the unit
\emph{Endgame} applies it twice, once to bound the height of
\m{\chainVolume{D}{1}^2/\chainVolume{D}{\theta_{\ell,\iota}}^2} for a fit family~$\theta$, and once for the
non-vanishing in~(b). The other input bounds each ratio
\m{\liftVolume{D}{1}{j}/\liftVolume{D}{\theta_{\ell,\iota}}{j}}; it is proved in the note \emph{The height of a ratio
of lift volumes}, which uses Lemmas \ref{l:idem} and~\ref{l:range} to describe the image of a single lift, and
Lemma~\ref{l:equiv} is used in the note \emph{Glued lattices, product forms and the integral comparison forms}.
Chapter~VIII uses Lemmas \ref{l:idem} and~\ref{l:range} in the proof that fit families exist.

\sect{Notes on the formal text}{s:notes}
\begin{enumerate}[label=\textup{\arabic*.},leftmargin=1.6em,itemsep=5pt]
\item \emph{The volume factor.} The formal text uses the library's volume factor of a linear map and its two
properties quoted in Definition~\ref{d:vol}: its square is \m{\det(u^{*}u)}, and it vanishes exactly when the
kernel is not zero. The library's definition was not read; the definition above is by these properties. The lift
volume and the volume of the chain are defined as squares of this real number, cast to~$\CC$.
\item \emph{The constant is explicit.} The formal proof gives $c$ by the formula at the end of the proof above. It
does not say that $c$ is unique, and the statement does not use more than its existence.
\item \emph{Two identifications made by Lean.} The proof applies Input~5 to the series of Input~2, for the
twist~$\vartheta$ and the row~$i$, and to the series of Input~1, for the tame character \m{\vartheta^{i+1}} and
the lift index \m{i+1}. That the terms of the first series are the Rallis terms with the parameters
\m{(\vartheta^{i+1},i+1)}, and that the Gram factor of the pair \m{O_D(\vartheta^{i+1},i+1)} is the one of
\m{O_D(\vartheta,1)}, is nowhere stated as a lemma: Lean accepts the application because the definitions unfold
to the same terms. Input~2 above is worded with these identifications made.
\item \emph{Arguments that the formal statements leave to Lean.} In the definition of the volume of the chain the
twist is written without its type, and in the definition of the lift volume the index~$j$ is; Lean infers that
the first is an $m$-fold twist and the second a lift index, from the maps they are given to. In
Lemma~\ref{l:equiv} the form~$y$ and in Lemmas \ref{l:equiv} and~\ref{l:range} the twist~$\vartheta$ are written
without types as well; that $\vartheta$ is a tame character, not necessarily a twist, was read off the definition
of the adjoint lift.
\item \emph{Hypotheses of Lemma~\ref{l:gram}.} The formal statement asks that both spaces be finite-dimensional;
$U'$ could be arbitrary once the adjoint exists, but the library's volume factor is defined for finite-dimensional
spaces.
\item \emph{What an assembly assembles.} Theorem~A carries no analysis of its own: the convergence of the Rallis
series, the triangle of lifts and its square laws, the rigidity of the labels behind Input~6 and the index
computation of Input~4 all lie in the inputs.
\end{enumerate}

\subhead{How this note was written.} From the dossier of the unit: the structured proofs of Theorem~A (eight
steps), of Lemma~\ref{l:gram} and of Lemmas \ref{l:idem} and~\ref{l:range}. I read the whole Lean file of these
eight declarations, and the Lean statements of the inputs in their files: the cost of squares and its two rules,
the inner product formula for one lift, the formula for the first lift of a row, the statement about the Gram
form of a weighted chain in a triangle together with the definition of a triangle, the index of the last level
and the norm of the chain at its level, the common Hecke element, and the two statements about good forms on
different spaces; also the Lean definitions of tame characters, twists, $m$-fold twists, forms of principal level,
good forms, rows, heights and fit families. I did not read the proofs of the inputs. The reader of the dossier
notes that no second reader had looked at it. After the dossier was frozen, the second reader of the write-up compared the readings of all statements of the unit
with the Lean text and sampled its structured proofs (her message to the writer of 26 November 2026); for this note
she confirms the two statements about volumes and that the last two lemmas (Lemmas \ref{l:equiv}
and~\ref{l:range}) take a tame character, as stated here. The proof of Theorem~A above follows the formal one step by step. The program check of the unit's dossier covered \UnitSteps{} leaf
steps of \UnitProofs{} proofs.

\subhead{Revision.} This text was revised after the report of an independent referee (26 November 2026), who
re-derived the proof of Theorem~A and compared the inputs with their Lean statements. The changes: Inputs 1 and~2
are no longer called instances of the Rallis inner product formula, and the references are those of the catalogue
with the second reader's verdict; the note says that the constant~$c$ does not vanish; the levels and the ring of
integers are introduced in the Setting; the implicit arguments of the last three lemmas are listed correctly.

\head{Appendix A. Concordance with the Lean artifact}
{\fontsize{8.6}{11.6}\selectfont\renewcommand{\leanfont}{\ttfamily\fontsize{7.7}{10}\selectfont\color{leancol}}
\setlength{\LTleft}{0pt}\setlength{\LTright}{\fill}\setlength{\tabcolsep}{0pt}
\begin{longtable}{@{}p{30mm}p{56mm}p{43mm}r@{}}
\textbf{In this note} & \textbf{Lean declaration} & \textbf{File} & \textbf{Lines}\\\midrule\endhead
\ConcordanceRows
\end{longtable}}

\subhead{Theorem A and the three definitions as they are written in the artifact.}
\quoted{volumeProduct}
\quoted{liftVolume}
\quoted{chainTheta}
\quoted{chainVolume}
In these texts {\leanfont D.MFoldTwist} is the type of the $m$-fold twists of~$D$, {\leanfont D.goodLift} the
adjoint lift restricted to the good forms, {\leanfont D.iotaAtLevel} the map \m{\chain_D(\vartheta)},
{\leanfont D.rows} the space \m{\Row(D)}, {\leanfont D.rowDim} its dimension, {\leanfont normDet} the volume
factor, and {\leanfont D.CostsSquares} the notion of Input~7.

\head{Notation}
\begingroup\fontsize{8.6}{10.9}\selectfont\raggedright
\setlength{\columnsep}{9mm}\setlength{\parskip}{0pt}
\begin{multicols}{2}
\nt{$P=(N_0,d)$, $D$}{parameters; admissible data for~$P$}
\nt{$E(D)$}{the CM field of the data}
\nt{$m$, $r$, $t(D)$, $x$}{number of lifts; rank of the rows; the step \m{5r}; the first scale}
\nt{$g(D)$, $N(D)$}{the two levels of the data}
\nt{$\sfG(D)$}{the finite group of the chosen primes}
\nt{$j$, $i$, $n$, $e$}{a lift index; a row index \m{<m}; integers \m{\ge0}}
\nt{$Y_i(D)$, $W_j(D)$}{the hermitian spaces of the chain; of the $j$-th single lift}
\nt{$V$, $V'$}{hermitian spaces of rank~$r$}
\nt{$\Mf_D(V)$, $\Mg_D(V)$}{forms of principal level on~$V$; the good ones}
\nt{$\Row(D)$, $n_0$}{the rows; their dimension}
\nt{$e_0(D)$, $\idem(\mathbf{c})$, $\mathbf{c}$}{the good idempotent; the idempotent of a tuple of labels; a tuple}
\nt{$\Hk_D$, $T$, $T_V$, $T_{\vartheta,i}$}{the Hecke algebra; an element; its action on~$\Mf_D(V)$; the element
of the proof}
\nt{$p$, $\Mf_D(V)[p]$}{a Hecke eigensystem; its eigenspace}
\nt{$\vartheta$, $Q(\vartheta)$}{a tame character or twist; its set of ramified primes}
\nt{$\mathfrak{m}_D(\vartheta,e)$, $\mathcal{O}$}{the ideal of chosen primes to the power~$e$ times
\m{Q(\vartheta)}; the ring of integers of \m{E(D)}}
\nt{$O_D(\vartheta,j)$, $\Ff(O_D(\vartheta,j))$}{the auxiliary pair; its space of algebraic forms}
\nt{$\alift_D(\vartheta,j)$}{the adjoint lift}
\nt{$\chain_D(\vartheta)$, $\mathcal{A}_{D,\vartheta}$}{the chain of $m$ lifts; its target}
\nt{$\chainTheta{D}{\vartheta}$}{the chain restricted to the rows}
\nt{$\vol(u)$, $u^{*}$}{the volume factor of a linear map~$u$; its adjoint}
\nt{$\liftVolume{D}{\vartheta}{j}$, $\chainVolume{D}{\vartheta}$}{the lift volume; the volume of the chain}
\nt{$L$}{in ``$L$-function'': the Dirichlet series attached to a character; in the remark after the inputs a theta
lift}
\nt{$\ell$, $\iota$, $\sigma$}{a prime; a field isomorphism from~$\CC$ to an algebraic closure of~$\mathbb{Q}_\ell$;
an automorphism of~$\CC$}
\nt{$\theta=(\theta_{\ell,\iota})$}{a family of $m$-fold twists}
\nt{$Z$, $h$, $H$, $\alpha$, $\eta_\ell$}{a family of complex numbers; the objects of a height bound}
\nt{$c$, $c'$, $c_0$, $\kappa$}{functions on the $m$-fold twists}
\nt{$A_j$, $C_i$, $G$, $G'$, $l$, $l'$, $\omega$}{the constants of the inputs}
\nt{$\delta_D(\vartheta)$, $I(\vartheta)$, $N_n(\vartheta)$}{the Gram factor; the index of the last level; the
weights of the triangle}
\nt{$R_D(\vartheta,j,V;y',y;\gamma)$, $\gamma$}{a Rallis term; a coset}
\nt{$c_F$, $\rho(\gamma)$, $R(\gamma)$}{in the remark after the inputs: the mass of a finite datum, a matrix
coefficient, a right translation}
\nt{$F_\vartheta(i,n)$, $a_\vartheta(i,n)$, $\jmath_i$, $u_\vartheta$}{the triangle: spaces, lifts, embeddings;
the weighted chain}
\nt{$S(\vartheta,i)$, $T(\vartheta)$, $S$}{endomorphisms of the rows; an endomorphism}
\nt{$U$, $U'$, $u$, $a$, $y$, $y'$, $z$}{inner product spaces, a linear map, a scalar, vectors or forms}
\nt{$A$, $A'$, $B$, $B'$}{constants and forms of Input~5}
\nt{$\varrho$, $\xi$, $\mathcal{N}$, $y_{\mathbf{c}}$}{a homomorphism of algebras; an element of the group algebra;
a subspace; a component}
\nt{$n'$, $n''$, $j'$}{further integers and a further lift index}
\end{multicols}
\endgroup

\end{document}
